\chapter{Introduction} \label{chap:intro}

Coastal ocean environments are highly dynamic systems, shaped by the interaction of atmospheric forcing, bathymetry, coastal geometry, circulation patterns and local mixing processes. These interactions generate spatially heterogeneous temperature fields, fronts, gradients and transient structures that are difficult to observe with sufficient spatial and temporal resolution. Accurate knowledge of these structures is essential for improving ocean forecasts and supporting scientific, environmental and operational decision-making in coastal regions \cite{wunsch2007stateestimation,dowd2014dataassimilation}.

Numerical ocean models provide continuous estimates of the ocean state, but their predictive skill depends strongly on the availability, quality and spatial distribution of observations. In situ measurements are often sparse, costly and logistically demanding, while remote sensing products are limited by atmospheric conditions, surface-only coverage and resolution constraints. Autonomous Underwater Vehicles (AUVs) offer a flexible alternative for targeted ocean sampling, as they can collect high-resolution measurements along adaptive trajectories and complement existing observing systems \cite{curtin2009aosn,hwang2019auvreview}. Nevertheless, AUV missions are constrained by endurance, navigation requirements, safety limitations, communication restrictions and operational costs. These constraints make trajectory planning a central component of efficient ocean observation.

Adaptive sampling strategies address this challenge by using model information to guide the selection of informative sampling locations. In model-driven approaches, ocean forecasts and their associated uncertainty can be used to define value maps that indicate where new observations are expected to be most useful \cite{lermusiaux2007adaptive,mendes2025oceanmodeldriven}. In uncertainty-driven planning, pointwise standard deviation maps can be used to prioritise regions where the model prediction is less certain, supporting the design of AUV trajectories that maximise information collection under operational constraints \cite{bernacchi2025closingloop,duarte2025geostatistical}.

However, uncertainty maps do not necessarily encode all relevant spatial structures present in the underlying ocean field. Temperature maps may reveal thermal regimes, gradients, fronts and boundary-like structures that provide complementary information about the organisation of the ocean state. This distinction is particularly relevant in coastal regions where different dynamical regimes may coexist within the same operational domain. In such cases, observations collected in one region may not be equally informative for another, and model updates or sampling decisions based only on uncertainty may fail to account for regime-dependent behaviour \cite{mendes2025oceanmodeldriven}. Previous work on compact ocean models and regime-based representations has shown that spatial patterns in oceanographic fields can be exploited to characterise recurrent environmental states and support robotic sampling decisions \cite{fossum2020compact,ge2023watermasses}. This motivates the use of temperature-derived descriptors as additional information for adaptive planning.

This dissertation investigates whether spatial descriptors extracted from ocean temperature fields can complement an uncertainty-driven adaptive planner for AUV missions. The work focuses on the Nazaré Canyon region, where complex bathymetry and coastal dynamics can generate distinct thermal structures and regime transitions. High-resolution temperature and uncertainty products derived from ocean model data and geostatistical prediction are used as the basis for the proposed pipeline. The approach compares a baseline planner driven by uncertainty maps with enriched planning formulations that incorporate descriptors related to thermal regimes, boundaries and gradients.


\section{Context and motivation} \label{sec:context_motivation}

The motivation for this dissertation arises from the FRESNEL experiment, which demonstrated the feasibility of coupling ocean model forecasts, geostatistical uncertainty estimation, targeted AUV sampling and subsequent data assimilation within a sample--assimilate--predict--direct loop \cite{mendes2025oceanmodeldriven}. In that framework, model-derived uncertainty maps were used to guide AUV trajectories towards regions where new observations were expected to be most informative. This established uncertainty-driven planning as a practical and operationally relevant baseline for model-driven robotic exploration.

However, the FRESNEL results also revealed an important limitation. Although assimilating AUV observations generally improved the statistical model predictions, the impact was not spatially uniform. In particular, observations collected in one part of the domain did not necessarily improve predictions in another part of the region. The XP2 case showed that assimilation could even degrade prediction skill when the assimilated observations were collected south of the Nazaré Canyon while the validation trajectory was located to the north. This behaviour was interpreted as a consequence of spatial mismatch and the presence of distinct dynamical regimes across the canyon region \cite{mendes2025oceanmodeldriven}.

This observation is central to the present work. If different parts of the operational domain exhibit distinct thermal or dynamical regimes, then an uncertainty-only reward map may not fully represent the structural organisation of the ocean state. A trajectory that maximises accumulated uncertainty may collect valuable observations locally, but it may not explicitly account for regime boundaries, transitions or representative regions that influence how informative those observations are for the broader domain.

Temperature fields provide a natural source of additional spatial information. Beyond their pointwise values, they contain gradients, fronts, boundary-like structures and recurrent spatial patterns that may reflect distinct oceanographic regimes. Compact representations of ocean fields have previously been used to identify dominant spatial states and support adaptive robotic sampling \cite{fossum2020compact,ge2023watermasses}. These ideas motivate the use of temperature-derived descriptors as complementary information for AUV planning.

Therefore, this dissertation investigates whether descriptors extracted from temperature-regime representations can complement the uncertainty-driven planning approach used in FRESNEL. The aim is not to replace uncertainty as the main planning input, but to evaluate whether adding regime-aware information can modify trajectory behaviour in a meaningful and operationally interpretable way.


\section{Problem statement} \label{sec:problem_statement}

The FRESNEL experiment demonstrated that uncertainty-driven AUV sampling can support short-term ocean prediction by coupling model forecasts, geostatistical uncertainty estimation, targeted robotic sampling and data assimilation within a sample--assimilate--predict--direct loop \cite{mendes2025oceanmodeldriven}. In this framework, standard deviation maps are used as reward maps, guiding AUV trajectories towards regions where the model prediction is less certain and where new observations are expected to be informative.

However, the FRESNEL results also revealed that the benefit of assimilated observations can be spatially dependent. In the Nazaré Canyon region, observations collected in one part of the operational domain did not necessarily improve predictions in another. In particular, the XP2 case showed that assimilation could degrade prediction skill when the assimilated observations were collected south of the canyon while the validation trajectory was located to the north \cite{mendes2025oceanmodeldriven}. This behaviour suggests that the operational region cannot always be treated as a single homogeneous domain and that distinct thermal or dynamical regimes may influence the usefulness of collected observations.

The baseline planning strategy used in this context is primarily uncertainty-driven. It is well suited for targeting regions of high model uncertainty, but it does not explicitly encode the spatial organisation of the predicted temperature field, such as gradients, fronts, regime transitions or representative regions. Therefore, a trajectory that maximises accumulated uncertainty may be locally informative, while still not explicitly accounting for the regime structure that may affect how observations relate to the broader domain.

The problem addressed in this dissertation is how to introduce regime-aware spatial information into an uncertainty-driven AUV planner without replacing its original uncertainty-based objective. Given predicted temperature maps and associated standard deviation maps over the Nazaré Canyon region, this work investigates how recurrent temperature regimes can be identified, how descriptors can be extracted from their representative patterns, and how those descriptors can be integrated into the planner reward formulation.

The central question is whether descriptor-enriched planning can produce trajectories that remain uncertainty-seeking while becoming more sensitive to thermal structure, regime transitions and spatial heterogeneity. This work does not aim to prove a direct improvement in data assimilation. Instead, it evaluates whether temperature-derived descriptors can provide complementary planning information that addresses a limitation observed in the FRESNEL mission and supports a more regime-aware interpretation of AUV trajectory planning.


\section{Objectives} \label{sec:objectives}

The main objective of this dissertation is to investigate whether spatial descriptors extracted from ocean temperature fields can complement an uncertainty-driven adaptive planner for AUV missions in the Nazaré Canyon region.

This objective is structured into the following specific goals:

\begin{itemize}
  \item \textbf{O1 --- Build a temperature-regime discovery pipeline:} develop and validate a pipeline capable of identifying spatial regimes or recurrent patterns in high-resolution ocean temperature maps over a defined region of interest.

  \item \textbf{O2 --- Extract planning-relevant descriptors:} derive spatial descriptors from the temperature-regime representation, including descriptors related to thermal gradients, boundary regions, regime transitions and representative spatial structures.

  \item \textbf{O3 --- Integrate descriptors with an uncertainty-driven planner:} incorporate the extracted descriptors into an existing AUV planner that uses standard deviation maps as its baseline reward input, while preserving the original uncertainty-driven planning logic.

  \item \textbf{O4 --- Compare baseline and descriptor-enriched planning formulations:} evaluate trajectories generated by a standard deviation-only baseline planner against enriched formulations that combine uncertainty with temperature-derived descriptors under the same grid, operational constraints and planning scenarios.

  \item \textbf{O5 --- Assess sensitivity and operational feasibility:} analyse the effect of descriptor weights, grid choices, single- and multi-AUV configurations, and computational cost on the resulting trajectories and planner behaviour.
\end{itemize}




\section{Scope and assumptions} \label{sec:scope_assumptions}

This dissertation focuses on the planning layer of model-driven AUV sampling. The work assumes that forecast products, uncertainty maps and bathymetric or operational masks are provided by an external modelling workflow. In particular, temperature predictions and pointwise standard deviation maps are treated as available inputs to the planner. The development of new ocean forecasting or data assimilation methods is outside the scope of this dissertation.

The thesis is centred on two coupled components. The first is the extraction of spatial regimes and descriptors from gridded ocean temperature fields. The second is the use of these descriptors to enrich an uncertainty-driven trajectory planner. The planner is not redesigned from first principles; instead, the work investigates how additional descriptor-based reward terms or incentives can be incorporated into an existing planning framework.

The study is restricted to gridded products over a selected region of interest in the Nazaré Canyon area. The final dataset is expected to use a common high-resolution grid containing, at minimum, temperature, standard deviation, latitude, longitude and bathymetry information. This common grid is assumed to be the canonical spatial domain for both regime discovery and planner evaluation. Details on the final grid, region of interest and available files are provided in Chapter~\ref{chap:problem_specification}.

The work does not address the mechanical design of AUV platforms, low-level guidance and control, acoustic communication modelling, simultaneous localisation and mapping, or the development of new sensors. It also does not attempt to prove that the proposed enriched planner is universally optimal. \textcolor{red}{Instead, the objective is to provide a controlled evaluation of whether temperature-derived descriptors can alter planner behaviour in meaningful ways when compared with an uncertainty-only baseline.}

Although the planner may support single- and multi-AUV configurations, the main focus is the integration and evaluation of the descriptor-enriched planning concept. Multi-AUV behaviour is therefore considered as an evaluation scenario rather than as the central contribution of the dissertation.
\section{Expected contributions} \label{sec:expected_contributions}

\textcolor{red}{The expected contributions of this dissertation are:}

\begin{itemize}
  \item A \textbf{structured pipeline for temperature-regime discovery} in high-resolution ocean temperature maps, designed to identify spatial patterns and regime-like structures relevant to adaptive sampling.

  \item A set of \textbf{temperature-derived spatial descriptors}, including descriptors associated with gradients, boundaries and regime transitions, suitable for representation on the planner grid.

  \item An \textbf{integration strategy for descriptor-enriched AUV planning}, where the baseline uncertainty-driven planner is complemented with additional descriptor-based terms without replacing the standard deviation map as the main uncertainty input.

  \item A \textbf{baseline versus enriched evaluation framework}, comparing uncertainty-only planning with descriptor-enriched formulations under controlled scenarios, including sensitivity analysis for descriptor weights.

  \item A \textbf{case-study evaluation in the Nazaré Canyon region}, using high-resolution temperature and uncertainty products to assess whether regime-aware descriptors can support more structured and oceanographically meaningful sampling trajectories.

  \item A \textbf{reproducible implementation and validation workflow}, including data/grid validation, planner interface generation, diagnostic plots and audit artefacts supporting the interpretation of the results.
\end{itemize}

\section{Dissertation structure \textcolor{red}{No final de tudo revejo}} \label{sec:struct}

In addition to this introductory chapter, the dissertation is organised as follows.

Chapter~\ref{chap:background} presents the background and fundamental concepts required for the thesis, including ocean modelling and forecast uncertainty, autonomous underwater vehicles, adaptive sampling, geostatistical prediction, spatial regimes in ocean fields and the path-planning formulation used by the baseline planner.

Chapter~\ref{chap:related_work} reviews the related work and state of the art in adaptive ocean sampling, model-driven robotic exploration, uncertainty-driven trajectory planning, compact representations and regime discovery in marine robotics. It concludes by identifying the research gap addressed by this dissertation.

Chapter~\ref{chap:problem_specification} defines the problem specification and study case in concrete terms. It introduces the Nazaré Canyon operational region, the available datasets and formats, the variables used in the thesis, the operational constraints, functional and non-functional requirements, and the evaluation metrics.

Chapter~\ref{chap:proposed_approach} presents the proposed approach, including the overall pipeline, temperature-regime discovery, descriptor extraction, the selected grid and region of interest, the baseline uncertainty-driven planner and the enriched planner formulations.

Chapter~\ref{chap:implementation} describes the implementation details, including data preprocessing, the Fossum-style regime discovery pipeline, planner interface generation, integration with Lucrezia's planner and reproducibility artefacts.

Chapter~\ref{chap:validation} describes the verification, validation and experimental design. It includes data and grid validation, baseline versus enriched planner comparison, sensitivity analysis for descriptor weights, single- and multi-AUV scenarios, and computational cost assessment.

Chapter~\ref{chap:results} presents and discusses the experimental results.

Finally, Chapter~\ref{chap:conclusions} summarises the main findings, discusses limitations and outlines directions for future work.


% =========================
% CHAPTER 2 — BACKGROUND AND FUNDAMENTALS
% File: chapter2.tex
% =========================

\chapter{Background and Fundamentals} \label{chap:background}

This chapter introduces the fundamental concepts required to understand the problem addressed in this dissertation. The focus is not on comparing specific methods from the literature, which is addressed in Chapter~\ref{chap:related_work}, but on defining the main technical and scientific building blocks used throughout the thesis. These include ocean modelling and forecast uncertainty, autonomous underwater vehicles and adaptive sampling, geostatistical prediction and uncertainty maps, spatial regime discovery in ocean fields, and path-planning formulations based on graph routing and prize-collecting vehicle routing problems.

\section{Ocean modelling and uncertainty} \label{sec:ocean_modelling_uncertainty}

Numerical ocean models provide spatially and temporally continuous estimates of the ocean state. They are used to represent variables such as temperature, salinity, currents and sea surface height, and they support a wide range of scientific and operational applications \cite{wunsch2007stateestimation,lermusiaux2007adaptive}. In coastal regions, however, ocean modelling is particularly challenging because the dynamics are influenced by bathymetry, coastline geometry, atmospheric forcing, tides, mesoscale circulation and local mixing processes. These mechanisms generate heterogeneous fields with fronts, gradients and transient structures that may evolve over short time scales \cite{mendes2025oceanmodeldriven}.
% TODO: update citation format for mendes2025oceanmodeldriven depending on the final publication status.

A model prediction can be understood as an estimate of the true ocean state. Let $x(\mathbf{s},t)$ denote an ocean variable at spatial location $\mathbf{s}$ and time $t$, such as temperature. A forecast model provides an estimate
\begin{equation}
    \hat{x}(\mathbf{s},t),
\end{equation}
which approximates the unknown true field. The difference between the forecast and the true state is affected by several sources of error, including imperfect initial conditions, unresolved physical processes, limited spatial resolution, uncertain forcing and sparse observations \cite{dowd2014dataassimilation}. As a result, model predictions should not be interpreted only as deterministic fields, but also in terms of their associated uncertainty.

Forecast uncertainty can be represented in different ways. In ensemble-based or simulation-based approaches, multiple possible realisations of the ocean field are generated. The variability among those realisations provides a quantitative indication of uncertainty \cite{duarte2025geostatistical}. A common pointwise measure is the standard deviation:
\begin{equation}
    \sigma(\mathbf{s},t) =
    \sqrt{
    \frac{1}{M-1}
    \sum_{m=1}^{M}
    \left(x_m(\mathbf{s},t)-\bar{x}(\mathbf{s},t)\right)^2
    },
\end{equation}
where $M$ is the number of realisations, $x_m(\mathbf{s},t)$ is the value of the $m$-th realisation at location $\mathbf{s}$ and time $t$, and $\bar{x}(\mathbf{s},t)$ is the ensemble mean. In this dissertation, standard deviation maps are treated as uncertainty maps that can guide adaptive sampling decisions.

Uncertainty is important for robotic sampling because it provides a model-derived indication of where new observations may be valuable. Regions with higher uncertainty are often interpreted as areas where additional data may reduce forecast error or improve the reliability of subsequent predictions \cite{bernacchi2025closingloop,duarte2025geostatistical}. However, uncertainty is not the only relevant property of an ocean field. Temperature maps may also contain spatial patterns, such as fronts, gradients and regime transitions, that are relevant for planning and interpretation. This motivates the combination of uncertainty information with descriptors extracted from the predicted temperature field.

\section{Autonomous underwater vehicles and adaptive sampling} \label{sec:auv_adaptive_sampling}

Autonomous Underwater Vehicles (AUVs) are robotic platforms capable of collecting oceanographic measurements along planned trajectories. Compared with fixed stations or ship-based surveys, AUVs provide mobility and flexibility, allowing observations to be collected over spatially extended regions and at higher resolution along the vehicle path \cite{curtin2009aosn,hwang2019auvreview}. Typical payloads may include sensors for temperature, salinity, pressure, dissolved oxygen, chlorophyll fluorescence, turbidity and other physical or biogeochemical variables, depending on the mission objectives \cite{sousa2012lauv}.

Despite these advantages, AUV missions are constrained by several operational factors. Vehicle endurance limits mission duration and distance travelled. Safety constraints restrict navigation in regions with shallow bathymetry, obstacles, intense traffic or other operational hazards. Communication underwater is intermittent and low-bandwidth, which limits the possibility of continuous replanning or supervision. Currents may also affect travel time and energy consumption \cite{aguiar2018trajectoryopt}. These constraints mean that an AUV cannot sample everywhere, and the selection of sampling locations becomes a planning problem.

Adaptive sampling refers to strategies in which the sampling plan is informed by prior knowledge, model predictions, uncertainty estimates or data collected during previous missions \cite{lermusiaux2007adaptive,hwang2019auvreview}. Instead of following a fixed pre-defined survey pattern, the vehicle trajectory is designed to prioritise locations that are expected to be informative. In model-driven adaptive sampling, the planning process uses model products, such as predicted fields and uncertainty maps, to identify valuable regions for observation \cite{mendes2025oceanmodeldriven}.

In uncertainty-driven adaptive sampling, the planner favours regions where the model is less certain. This is particularly useful when the objective is to improve model predictions or collect data in areas where the forecast is expected to be least reliable. In practice, the uncertainty map can be converted into a reward or value map, and the AUV trajectory is selected to maximise the accumulated reward under mission constraints \cite{bernacchi2025closingloop,duarte2025geostatistical}. This principle underlies the baseline planner considered in this dissertation.

Adaptive sampling can also be extended beyond uncertainty-only criteria. If additional descriptors are available, such as thermal gradients, regime boundaries or representative regions, these can be incorporated into the planning objective. \textcolor{red}{The aim is not necessarily to replace uncertainty as the main driver, but to complement it with information about the spatial structure of the ocean field \cite{fossum2020compact,fossum2019phytoplankton}.}

\section{Geostatistical prediction and uncertainty maps} \label{sec:geostatistical_prediction}

Geostatistical methods provide a statistical framework for modelling spatially distributed variables. In ocean applications, they can be used to generate short-term predictions and associated uncertainty estimates from a combination of prior model fields and observational data \cite{cressie2011spatiotemporal,duarte2025geostatistical}. Unlike full numerical ocean models, geostatistical approaches can be computationally efficient and suitable for operational workflows where predictions and uncertainty maps need to be generated rapidly \cite{mendes2025oceanmodeldriven}.

A geostatistical prediction workflow typically relies on a covariance or variogram model that describes how values of the variable of interest are correlated in space and time. Given conditioning information, such as previous model fields or observations, stochastic simulation can be used to generate multiple possible realisations of the target field. The median or mean of these realisations can be used as a predicted field, while the pointwise standard deviation across the realisations provides an uncertainty map \cite{duarte2025geostatistical,mendes2025oceanmodeldriven}.

Let $\{x_1(\mathbf{s}), x_2(\mathbf{s}), \dots, x_M(\mathbf{s})\}$ be an ensemble of simulated temperature fields at location $\mathbf{s}$. A predicted temperature field may be obtained from the ensemble median:
\begin{equation}
    \hat{x}_{\mathrm{med}}(\mathbf{s}) =
    \mathrm{median}\left(x_1(\mathbf{s}), x_2(\mathbf{s}), \dots, x_M(\mathbf{s})\right),
\end{equation}
while the corresponding uncertainty field may be represented by the standard deviation $\sigma(\mathbf{s})$. These two products, predicted temperature and uncertainty, are central to this dissertation.

In the planning context, the predicted temperature field provides information about the expected spatial structure of the ocean state. The uncertainty map provides information about where the model is less confident. These two maps have different meanings: the temperature map describes the predicted value of the physical field, while the standard deviation map describes the variability or uncertainty of the prediction. This distinction is important because the proposed approach derives regime descriptors from temperature fields and combines them with uncertainty maps for planning.

Operationally, geostatistical uncertainty maps can be used to define regions of high expected sampling value. If a planner receives a map $\sigma(\mathbf{s})$, locations with larger values may be assigned higher reward. In this way, the uncertainty field becomes an actionable input to a trajectory planner \cite{bernacchi2025closingloop,duarte2025geostatistical}. The enriched approach explored in this thesis builds on this idea by adding descriptor-based information derived from the corresponding temperature field.

\section{Spatial regimes and pattern discovery in ocean fields} \label{sec:spatial_regimes}

Ocean fields are rarely spatially uniform. Temperature, salinity and related variables often contain coherent regions, gradients and transitions that reflect underlying physical processes. In coastal environments, these structures may result from the interaction between bathymetry, coastline geometry, atmospheric forcing, circulation patterns, upwelling and relaxation events, and local mixing processes. In canyon regions, topographic forcing can further contribute to distinct dynamical behaviour across different parts of the domain \cite{mendes2025oceanmodeldriven}.

In this dissertation, a temperature regime refers primarily to a recurrent spatial pattern of the full temperature field, rather than to an individual grid-cell label. This definition follows compact-model approaches in marine robotics, where oceanographic maps with similar magnitudes and spatial patterns are grouped into dominant classes or representative states \cite{fossum2020compact}. Once a daily temperature map is assigned to a regime class, descriptors are extracted from the corresponding class prototype to represent local structures such as gradients, boundaries, cold and warm regions, and representative zones.

This distinction is important for the proposed methodology. The regime discovery step classifies full temperature maps according to their spatial similarity, producing a compact representation of recurrent thermal patterns. The descriptor extraction step then analyses the representative prototype of each class in order to identify spatial features that may be relevant for planning. Therefore, the descriptors used in the planner do not correspond to independent cell-wise regime labels; instead, they are derived from the spatial organisation of the class prototype associated with the predicted temperature field.

Pattern discovery aims to identify recurrent structures from data. In gridded temperature fields, this may involve clustering days or maps according to their spatial similarity, detecting recurring patterns, or extracting descriptors that summarise relevant spatial information. Compact representations and clustering-based approaches have previously been used to summarise oceanographic fields and support adaptive sampling decisions \cite{fossum2020compact,ge2023watermasses}. These approaches are particularly relevant when the objective is to reduce a high-dimensional environmental field to a smaller set of interpretable states or patterns that can inform robotic sampling.

A simple example of a spatial descriptor is the gradient magnitude of a temperature field:
\begin{equation}
G(\mathbf{s}) = \left| \nabla T_{\mathrm{norm}}(\mathbf{s}) \right|,
\end{equation}
where $T_{\mathrm{norm}}(\mathbf{s})$ denotes the normalised temperature field at location $\mathbf{s}$. High values of $G(\mathbf{s})$ indicate regions where temperature changes rapidly in space, which may correspond to fronts, transition zones or sharp thermal gradients. In this dissertation, such descriptors are computed on normalised temperature fields or class prototypes, so that the emphasis is placed on spatial structure rather than on absolute temperature magnitude alone.

Similarly, boundary-related descriptors can be used to represent transitions between different spatial structures within a class prototype. A boundary descriptor may assign higher values to cells located near transitions between cold and warm regions, between representative and non-representative areas, or near spatial discontinuities identified from the prototype. Such descriptors are motivated by the fact that fronts and boundaries between water masses can be informative targets for adaptive sampling, particularly when the objective is to characterise spatial heterogeneity or transitions between oceanographic regimes \cite{ge2023watermasses}.

The regime discovery component of this thesis is therefore inspired by compact representations of ocean fields, where high-dimensional spatial maps are transformed into lower-dimensional classes or prototypes that summarise recurrent structures \cite{fossum2020compact}. In this work, these classes are not used as a final classification product only. Instead, their prototypes are used to derive planning-relevant descriptors that provide additional information to the uncertainty-driven planner. The central idea is that uncertainty maps indicate where the model prediction is uncertain, while regime descriptors indicate where the temperature field is structurally informative.

This distinction underpins the descriptor-enriched planning approach developed in this dissertation. The baseline planner remains driven by model-derived uncertainty maps, while the proposed descriptors provide complementary regime-aware information related to gradients, boundaries, representative regions and spatial heterogeneity.


\section{Path planning and PC-VRP formulation} \label{sec:pcvrp}

Once a value map is defined, the planner must convert it into one or more feasible AUV trajectories. This requires balancing the reward obtained by visiting informative locations against the cost of travelling between them. The problem can be represented as a graph, where nodes correspond to candidate sampling locations and edges correspond to travel costs. Each node has an associated reward, derived for example from uncertainty, descriptor values or a combination of both \cite{bernacchi2025closingloop}.

Let $V$ be a set of candidate nodes and let $r_j$ denote the reward associated with node $j \in V$. Let $c_{ij}$ denote the cost of travelling between nodes $i$ and $j$. A planner seeks a route or set of routes that maximises collected reward while respecting operational constraints such as endurance, start and end locations, obstacle masks and admissible regions. In a simplified form, the objective can be written as:
\begin{equation}
    \max_{\mathcal{R}} \sum_{j \in \mathcal{R}} r_j - \lambda \sum_{(i,j) \in \mathcal{R}} c_{ij},
\end{equation}
where $\mathcal{R}$ represents the selected route and $\lambda$ controls the trade-off between collected reward and travel cost.

This type of formulation is related to prize-collecting vehicle routing problems (PC-VRP). In a PC-VRP, not all candidate nodes need to be visited. Instead, the planner selects a subset of valuable nodes and constructs feasible routes that maximise reward under cost or time constraints \cite{bernacchi2025closingloop}. This is appropriate for AUV adaptive sampling because the vehicle cannot sample every cell of the domain and must prioritise locations according to the objective map.

In an uncertainty-driven planner, the node reward can be derived directly from the standard deviation map:
\begin{equation}
r_j^{\mathrm{base}} = \sigma_{\mathrm{norm}}(\mathbf{s}_j),
\end{equation}
where $\mathbf{s}*j$ is the spatial location of node $j$ and $\sigma*{\mathrm{norm}}$ is the normalised standard deviation map. This corresponds to the uncertainty-only baseline, in which higher reward is assigned to locations where the model prediction is less certain.

In a descriptor-enriched planner, the reward can be extended by combining the uncertainty map with a normalised descriptor map:
\begin{equation}
r_j^{\mathrm{enriched}} =
(1-\alpha),\sigma_{\mathrm{norm}}(\mathbf{s}*j)
+
\alpha,D*{\mathrm{norm}}(\mathbf{s}*j),
\end{equation}
where $D*{\mathrm{norm}}(\mathbf{s}_j)$ represents a normalised temperature-derived descriptor, such as a boundary score, gradient-related descriptor, representative-region map or other regime-aware spatial descriptor. The parameter $\alpha \in [0,1]$ controls the relative contribution of the descriptor term. The case $\alpha=0$ corresponds to the uncertainty-only baseline, while larger values of $\alpha$ increase the influence of the descriptor.

For multi-AUV experiments, vehicle-specific reward maps may also be considered. In this case, different vehicles can receive different descriptor incentives while sharing the same uncertainty map:
\begin{equation}
r^{(1)}*j =
w*{\sigma},\sigma_{\mathrm{norm}}(\mathbf{s}_j)
+
w_A,R_A(\mathbf{s}_j),
\end{equation}

\begin{equation}
r^{(2)}*j =
w*{\sigma},\sigma_{\mathrm{norm}}(\mathbf{s}_j)
+
w_B,R_B(\mathbf{s}_j),
\end{equation}
where $R_A$ and $R_B$ represent complementary region-aware descriptor maps assigned to different vehicles. This formulation allows the planner to remain primarily uncertainty-driven while testing whether regime-aware descriptors can influence trajectory structure and fleet-level spatial allocation.

The exact reward formulation used in each experiment depends on the descriptor selected, the normalisation procedure and the single- or multi-AUV configuration. However, the fundamental planning problem remains the same: select feasible AUV trajectories that collect high-value information under operational constraints. Chapter~\ref{chap:proposed_approach} describes how this dissertation builds on the baseline uncertainty-driven planner and extends it with temperature-derived descriptors.


\section{Summary} \label{sec:background_summary}

This chapter introduced the main concepts required for the remainder of the dissertation. Ocean models provide predicted fields and uncertainty estimates, while AUVs offer a flexible platform for targeted data collection under operational constraints. Geostatistical prediction provides a practical way to generate temperature and uncertainty maps for adaptive sampling. Temperature fields contain spatial structures such as gradients, recurrent regimes and boundaries, which can be represented through compact patterns and prototype-derived descriptors. Finally, path planning can be formulated as a graph-based reward maximisation problem, related to PC-VRP, where the planner balances information value against travel cost and feasibility constraints.

The following chapter reviews related work on adaptive ocean sampling, model-driven robotic exploration, uncertainty-driven planning, compact representations and regime-aware sampling, positioning the contribution of this dissertation within the existing literature.


\textcolor{red}{FIZ ATÉ AQUI}



% =========================
% CHAPTER 3 — RELATED WORK / STATE OF THE ART
% File: chapter3.tex
% =========================

\chapter{Related Work and the State of the Art} \label{chap:related_work}

This chapter reviews the literature most relevant to the integration of temperature-derived spatial descriptors into uncertainty-driven AUV planning. The aim is not to provide a broad review of all marine robotics or machine-learning approaches, but to position the dissertation within the specific intersection of adaptive ocean sampling, model-driven robotic exploration, uncertainty-based trajectory planning, compact representations of ocean fields, and regime-aware sampling.

The chapter is organised as follows. Section~\ref{sec:rw_adaptive_sampling_auvs} reviews adaptive sampling with AUVs. Section~\ref{sec:rw_model_driven_sampling} discusses model-driven ocean sampling, where ocean model products guide robotic decisions. Section~\ref{sec:rw_uncertainty_planning} focuses on uncertainty-driven trajectory planning. Section~\ref{sec:rw_clustering_compact_models} reviews clustering and compact models for marine robotics. Section~\ref{sec:rw_regime_spatial_mismatch} discusses regime-aware planning and spatial mismatch. Finally, Section~\ref{sec:rw_gap} identifies the research gap addressed by this dissertation.

% ==========================================================
\section{Adaptive sampling with AUVs} \label{sec:rw_adaptive_sampling_auvs}

Adaptive sampling refers to robotic sampling strategies in which the sampling plan is adjusted according to prior knowledge, model predictions, uncertainty estimates, or observations collected during the mission. Instead of following a fixed survey pattern, the vehicle selects sampling locations that are expected to provide higher scientific or operational value. This paradigm is particularly relevant in coastal ocean environments, where physical processes evolve rapidly and where fixed observing systems often cannot provide sufficient spatial coverage \parencite{hwang2019auvreview,lermusiaux2007adaptive}.

AUVs are well suited to adaptive sampling because they can collect high-resolution measurements along spatially extended trajectories while operating with a relatively low logistical footprint. Early and review work in marine robotics has shown that AUVs can support several adaptive sampling objectives, including mapping, feature tracking, front detection, plume monitoring, and uncertainty reduction \parencite{curtin2009aosn,hwang2019auvreview}. In these applications, the main challenge is not only to collect data, but to decide where data should be collected under constraints on time, energy, safety and navigation.

The literature distinguishes between several types of adaptive sampling objectives. Some approaches aim to reduce global uncertainty in a field, while others focus on tracking specific oceanographic features, such as fronts or biological patches. For example, coordinated robotic sampling has been used to track dynamic oceanographic features by combining model predictions, vehicle trajectories and feature-specific objectives \parencite{das2012coordinated}. Similarly, information-driven robotic sampling has been applied to coastal ocean environments where sampling decisions are guided by expected information gain or by the need to characterise spatially variable fields \parencite{fossum2018informationdriven,fossum2019phytoplankton}.

A common limitation of adaptive sampling strategies is the trade-off between scientific value and operational feasibility. A sampling location may be informative but unreachable within the mission duration, unsafe due to bathymetry or operational constraints, or redundant with other measurements. This motivates planning formulations that explicitly include both sampling reward and travel cost. These ideas are central to the uncertainty-driven and PC-VRP-based planning framework used as a baseline in this dissertation \parencite{bernacchi2025closingloop}.

For this dissertation, adaptive sampling provides the general robotic context: AUV trajectories are not treated as arbitrary paths, but as decisions driven by spatial value maps. For this dissertation, adaptive sampling provides the general robotic context: AUV trajectories are treated as decisions driven by spatial value maps. The thesis builds on this idea by investigating whether the uncertainty-based value maps used in model-driven planning can be complemented with descriptors extracted from temperature fields.

% ==========================================================
\section{Model-driven ocean sampling} \label{sec:rw_model_driven_sampling}

Model-driven ocean sampling uses numerical or statistical ocean model products to guide the planning of robotic observations. In this setting, the model provides a prediction of the ocean state and, when available, an estimate of the uncertainty associated with that prediction. The planner then uses this information to select sampling locations that are expected to improve knowledge of the field or support subsequent forecast updates \parencite{lermusiaux2007adaptive,curtin2009aosn}.

This approach is motivated by the fact that ocean models and observations are complementary. Models provide spatially and temporally continuous estimates of the ocean state, but their predictions are limited by uncertainty in initial conditions, forcing, parameterisations and unresolved processes. Observations provide direct information about the ocean, but are sparse and expensive to collect. Adaptive sampling aims to use robotic mobility to collect observations where they are most useful for constraining or evaluating the model \parencite{dowd2014dataassimilation,wunsch2007stateestimation}.

Recent work has demonstrated the feasibility of coupling ocean model predictions, uncertainty generation, robotic trajectory planning and data assimilation in real-world deployments. In such closed-loop frameworks, a model forecast is generated, an uncertainty map is derived, AUVs are directed to informative regions, and the collected measurements are assimilated to improve subsequent predictions \parencite{bernacchi2025closingloop,mendes2025oceanmodeldriven}. This model-driven view is directly relevant to the present dissertation, because the baseline planner relies on uncertainty fields derived from forecast products.

A practical challenge in operational model-driven sampling is computational cost. Full numerical ocean models can be expensive to run at high resolution and may not be suitable for rapid mission replanning. As an alternative, geostatistical or surrogate approaches can generate short-term predictions and uncertainty maps at lower computational cost, while still being conditioned on prior model information and available observations \parencite{duarte2025geostatistical,mendes2025oceanmodeldriven}. Such products can be transformed into planner-ready inputs, including predicted temperature maps and pointwise standard deviation maps.

However, model-driven sampling also inherits the assumptions and limitations of the underlying model. If the model or surrogate representation does not capture local heterogeneity, then the derived uncertainty and planning objectives may be incomplete. This is particularly important in complex coastal environments, where different regions of the domain may evolve according to distinct dynamical regimes. This limitation motivates the need to complement model-derived uncertainty with descriptors that encode spatial structure in the predicted field.

% ==========================================================
\section{Uncertainty-driven trajectory planning} \label{sec:rw_uncertainty_planning}

Uncertainty-driven trajectory planning is based on the idea that regions of high model uncertainty are valuable targets for observation. In this framework, a pointwise uncertainty field, such as a standard deviation map, is converted into a reward or value map. The planner then seeks a feasible trajectory that accumulates high reward while respecting operational constraints such as endurance, admissible regions, start and end points, and obstacle masks \parencite{bernacchi2025closingloop,duarte2025geostatistical}.

This approach has several advantages. First, it provides a direct connection between model uncertainty and sampling decisions. Second, it is interpretable: high-reward areas correspond to regions where the model is less certain. Third, it can be integrated with routing formulations that explicitly account for travel cost and mission constraints. In real-world AUV deployments, this makes uncertainty maps attractive as practical planning inputs \parencite{bernacchi2025closingloop}.

In the planner used as baseline in this dissertation, the uncertainty map is transformed into a set of candidate sampling locations. The resulting planning problem is related to a prize-collecting vehicle routing problem, where the vehicle is not required to visit every candidate node, but instead selects a subset that maximises collected reward under a travel budget \parencite{bernacchi2025closingloop}. Similar principles appear in broader vehicle routing formulations, where route feasibility, cost and reward must be balanced \parencite{toth2014vehicle,vidal2013hybrid}.

Uncertainty-driven planning has also been studied in relation to informative path planning and value-of-information approaches. These methods evaluate sampling actions according to their expected ability to reduce uncertainty or improve decision quality \parencite{eidsvik2015value,frolov2014grid}. In ocean sampling, this is particularly relevant because measurements are costly and the vehicle cannot cover the entire domain. The planner must therefore select a limited subset of locations expected to provide high information value.

Despite these advantages, uncertainty-driven planning has an important limitation: uncertainty is not equivalent to oceanographic structure. A region with high standard deviation may be uncertain but not necessarily located near an important front, boundary or regime transition. Conversely, a region with moderate uncertainty may still be scientifically relevant if it lies at the interface between distinct thermal regimes. Therefore, planning based only on uncertainty may overlook spatial structures that are relevant for understanding the ocean state.

This limitation motivates the enriched planning formulation explored in this dissertation. The baseline planner remains uncertainty-driven, but additional descriptors derived from the temperature field are considered as complementary information. In this way, the planner can be evaluated not only in terms of uncertainty reward, but also in relation to gradients, boundaries and regime-aware sampling behaviour.

% ==========================================================
\section{Clustering and compact models in marine robotics} \label{sec:rw_clustering_compact_models}

Marine robotic sampling often deals with high-dimensional spatial fields. Temperature maps, salinity fields, CTD profiles or satellite images can contain thousands of spatial values, making direct comparison and planning difficult. Compact representations address this problem by transforming high-dimensional observations into lower-dimensional features that preserve relevant structure. These representations can then be used for clustering, classification, anomaly detection or planning support.

Clustering methods are useful when the objective is to discover recurrent patterns without predefined labels. In oceanographic fields, clusters may correspond to recurring spatial structures, temperature regimes, water-mass configurations or feature states. Classification methods, in contrast, assign observations to predefined classes. Segmentation methods, in contrast, assign labels over a spatial domain and may reveal local boundaries between regions. Although this dissertation does not treat regime discovery primarily as a cell-wise segmentation problem, this distinction is useful because prototype-derived descriptors can still represent local structures such as boundaries and transitions. These concepts are relevant to this dissertation because the proposed pipeline seeks to extract spatial regime information from temperature maps and convert it into planning descriptors.

Compact models have been used in marine robotics to support adaptive sampling. \textcite{fossum2020compact} proposed compact representations based on dictionary learning and sparse coding, where high-dimensional ocean observations are encoded as sparse combinations of learned atoms. The sparse representation can then be used to identify recurring patterns and support adaptive sampling decisions. This is relevant to the present work because it demonstrates how spatial ocean fields can be converted into compact descriptors that preserve meaningful structure for planning.

Dictionary learning is particularly attractive because it can capture recurring spatial patterns while reducing dimensionality. Given a set of training samples, the method learns a dictionary of basis elements and represents each sample as a sparse combination of those elements. The resulting sparse codes can be clustered to identify classes or regimes. In the context of this dissertation, similar ideas motivate the use of temperature-map representations to identify spatial regimes and derive descriptors such as boundaries and gradients.

Other work has approached regime identification through probabilistic water-mass classification. For example, \textcite{ge2023watermasses} considered adaptive sampling for classification of water masses in three-dimensional environments. Such approaches are useful because they provide not only hard labels, but also uncertainty in the classification. This is especially important near boundaries between regimes, where the most informative sampling locations may lie.

However, many clustering or compact-model approaches are not directly integrated with operational AUV planners. They may identify patterns or regimes, but they do not necessarily show how these outputs should modify a trajectory planner that already uses uncertainty maps. This is the specific gap addressed by this dissertation: the objective is not only to identify regimes, but to transform regime information into descriptors that can be combined with an existing uncertainty-driven planner.

% ==========================================================
\section{Regime-aware planning and spatial mismatch} \label{sec:rw_regime_spatial_mismatch}

Regime-aware planning is motivated by the observation that coastal ocean fields are often heterogeneous. Distinct regions may differ in temperature structure, stratification, circulation, bathymetric influence or temporal evolution. In such conditions, observations collected in one part of the domain may not be equally informative for another. A planner that ignores these differences may select trajectories that maximise uncertainty reward, but fail to discriminate between regimes or to sample relevant transition zones.

The importance of this issue is illustrated by recent model-driven robotic exploration in the Nazaré Canyon region. In that experiment, targeted AUV observations generally improved geostatistical temperature predictions after assimilation. However, the improvement was not spatially uniform. For some trajectories, assimilation improved the forecast, whereas for others it produced weaker or even degraded performance. In particular, a spatial mismatch occurred when observations collected south of the Nazaré Canyon were used to evaluate or update predictions for a vehicle operating north of the canyon \parencite{mendes2025oceanmodeldriven}.
% TODO: update citation format for this reference depending on the final publication status of the manuscript.

This result is important because it indicates that the domain cannot always be treated as a single homogeneous region. The Nazaré Canyon can act as a transition zone between distinct oceanographic regimes, with circulation and thermal structure influenced by canyon topography and shelf-slope interactions. Under such conditions, sampling decisions and model updates may need to consider where observations are collected in regime terms, not only where the uncertainty is high \parencite{mendes2025oceanmodeldriven}.

This observation connects directly to the role of descriptors such as thermal gradients, boundaries and regime labels. Gradients can indicate where the field changes rapidly. Boundary descriptors can highlight transition zones between different regimes. Regime classes can indicate which full-field spatial patterns are recurrent, while descriptors extracted from their prototypes can highlight local structures such as boundaries, gradients and representative regions. When incorporated into a planner, these descriptors may encourage trajectories that sample relevant transitions, cover contrasting regions or avoid over-concentrating measurements within a single regime.

Previous work on feature-driven sampling, water-mass classification and compact models provides partial building blocks for this idea \parencite{das2012coordinated,fossum2020compact,ge2023watermasses}. However, there remains a gap between identifying regimes and using them as operational planning incentives in an uncertainty-driven AUV planner. This dissertation addresses that gap by evaluating descriptor-enriched planning formulations in which uncertainty remains the baseline objective, but temperature-derived regime information is added as complementary value.

% ==========================================================
\section{Research gap addressed by this thesis} \label{sec:rw_gap}

The reviewed literature provides several important foundations for this dissertation. Adaptive sampling with AUVs demonstrates the value of robotic mobility for collecting informative ocean observations \parencite{hwang2019auvreview,curtin2009aosn}. Model-driven ocean sampling shows how numerical or statistical forecasts can guide robotic missions \parencite{lermusiaux2007adaptive,mendes2025oceanmodeldriven}. Uncertainty-driven trajectory planning provides a practical way to convert standard deviation maps into reward fields for constrained route generation \parencite{bernacchi2025closingloop,duarte2025geostatistical}. Compact models and clustering demonstrate that spatial ocean fields can be summarised into regimes or lower-dimensional representations \parencite{fossum2020compact,ge2023watermasses}.

However, these components are not usually integrated into a single pipeline in which temperature-derived regime descriptors are explicitly used to enrich an uncertainty-driven AUV planner. In particular, the literature does not fully address how descriptors such as gradient magnitude, regime boundaries or representative regions should be combined with standard deviation maps in a planner that must still satisfy operational constraints. This is especially relevant in heterogeneous coastal environments, where spatial mismatch between regimes may limit the effectiveness of uncertainty-only sampling or assimilation.

The research gap addressed by this thesis can therefore be stated as follows: existing uncertainty-driven AUV planners can exploit standard deviation maps to target uncertain regions, but they do not explicitly account for the spatial organisation of the predicted temperature field. As a result, they may overlook thermal boundaries, regime transitions or structurally informative regions that are relevant for oceanographic interpretation and for avoiding regime-dependent sampling effects.

This dissertation addresses this gap by proposing and evaluating a descriptor-enriched planning pipeline. The pipeline first identifies spatial regimes and descriptors from high-resolution temperature maps. It then integrates these descriptors with a baseline uncertainty-driven planner. The resulting enriched formulations are compared against the baseline planner to assess whether temperature-derived descriptors can alter trajectory behaviour in meaningful and operationally feasible ways.

The contribution is therefore positioned at the interface between ocean modelling, marine robotics and spatial pattern discovery. It does not aim to replace numerical forecasting, geostatistical prediction or data assimilation. Instead, it investigates how additional structure extracted from predicted temperature fields can be used to improve the information available to an adaptive AUV planner.

% ==========================================================
\section{Summary} \label{sec:related_work_summary}

This chapter reviewed the related work required to position the dissertation. Adaptive sampling with AUVs provides the robotic context for targeted data collection. Model-driven ocean sampling links forecasts and uncertainty products to robotic decision-making. Uncertainty-driven trajectory planning shows how standard deviation maps can be converted into planner rewards. Clustering and compact models demonstrate how ocean fields can be represented as regimes or spatial patterns. Finally, recent evidence of spatial mismatch in heterogeneous coastal environments motivates regime-aware planning.

The gap addressed by this dissertation lies in the limited integration of temperature-derived spatial descriptors with existing uncertainty-driven AUV planners. The following chapter defines the problem specification and study case in concrete terms, including the Nazaré Canyon region, the available datasets, the variables used, operational constraints and evaluation metrics.


% =========================
% CHAPTER 4 — PROBLEM SPECIFICATION AND STUDY CASE
% File: chapter4.tex
% =========================

\chapter{Problem Specification and Study Case} \label{chap:problem_specification}

This chapter defines the operational problem, the study area, the data products and the evaluation quantities used in this dissertation. The central question is whether temperature-derived, regime-aware descriptors can complement an uncertainty-driven AUV planner in the Nazar\'e Canyon region. The baseline planner remains driven by the standard deviation map, hereafter denoted STD. The descriptors are evaluated as additional spatial information, not as a replacement for uncertainty and not as direct evidence of improved data assimilation.

\section{Operational motivation}

The motivation comes from the FRESNEL model-driven sampling workflow, in which model forecasts, geostatistical uncertainty maps, AUV sampling and data assimilation were combined in an operational loop. In that context, STD maps provided a practical uncertainty-driven reward for planning. However, the FRESNEL experience also highlighted that the usefulness of observations can be spatially dependent: samples collected in one part of the domain do not necessarily constrain errors elsewhere, especially near strong bathymetric and thermal structure around the Nazar\'e Canyon.

This dissertation therefore studies a complementary planning layer. Instead of changing the planner into a full data-assimilation system, it uses recurring temperature regimes and prototype-derived descriptors to reshape reward maps before they are passed to the existing Lucrezia planner. The resulting experiments assess trajectory behaviour, STD retention, descriptor sampling and operational feasibility under controlled comparisons.

\section{Study Area and Canonical Domain}

The study area is the FRESNEL/X490 region around the Nazar\'e Canyon, off the west coast of Portugal. The canonical historical dataset uses the same region of interest throughout the pipeline, which avoids comparing trajectories produced on different masks or grids. The input preparation starts from a larger CMEMS-derived surface temperature field, interpolates it to the high-resolution grid used in the local workflow and then extracts the X490 region used for the downstream regime analysis and planner inputs.

\begin{figure}[htbp]
    \centering
    \includegraphics[width=0.95\linewidth]{figures/results/cmems_to_x490_assets_contact_sheet.png}
    \caption{Data preparation from a larger CMEMS-derived surface temperature field to the interpolated grid and the final X490 region of interest around the Nazar\'e Canyon. The same ROI, valid mask and colour scale are used to support consistent visual comparison of the 370 historical surface temperature maps.}
    \label{fig:cmems_to_x490}
\end{figure}

The Step00 dataset audit reports 370 daily surface temperature maps from 2023-10-28 to 2024-10-31. The cropped temperature stack has shape $370 \times 72 \times 117$. The canonical common mask contains 8004 valid cells, corresponding to approximately 95.01\% of the ROI grid. The documented coordinate system is WGS84 / UTM zone 29N, with the requested X490 bounds spanning approximately 463--490 km in the east-west coordinate and 4376--4397 km in the north-south coordinate. The actual audited ROI bounds are close to these requested limits and cover approximately 557 km$^2$.

\begin{table}[htbp]
    \centering
    \small
    \caption{Audited canonical data used by the regime-inference and planning pipeline.}
    \label{tab:canonical_dataset}
    \begin{tabular}{p{0.32\linewidth}p{0.55\linewidth}}
        \hline
        Item & Audited value \\
        \hline
        Source variable & Surface temperature, \texttt{thetao} \\
        Number of daily maps & 370 \\
        Date range & 2023-10-28 to 2024-10-31 \\
        ROI & FRESNEL / X490 Nazar\'e Canyon ROI \\
        Cropped stack shape & $370 \times 72 \times 117$ \\
        Common valid cells & 8004 \\
        Valid fraction & Approximately 95.01\% \\
        Projection & WGS84 / UTM zone 29N, coordinates in km \\
        Approximate ROI area & 557 km$^2$ \\
        \hline
    \end{tabular}
\end{table}

\section{Planning-Day Products}

The planning-day stage uses two model-derived maps. The first is the predicted surface temperature field, denoted TEMPpred$(t)$, which is used to assign the planning day to one of the previously discovered full-field regime classes. The second is the STD map, denoted $\sigma(t)$, which remains the uncertainty-driven baseline reward. TEMPpred and STD are therefore used for different purposes: TEMPpred selects the appropriate prototype descriptor library, whereas STD defines the baseline information map for the planner.

The evaluated planning cases are associated with three representative dates used in the corrected Step12 experiments:

\begin{table}[htbp]
    \centering
    \small
    \caption{Planning cases used in the corrected planner evaluation.}
    \label{tab:planning_cases}
    \begin{tabular}{p{0.28\linewidth}p{0.18\linewidth}p{0.20\linewidth}p{0.20\linewidth}}
        \hline
        Case & Date & Assigned class & Main use \\
        \hline
        \texttt{C01\_representative} & 2024-08-24 & C01 & Single-AUV and multi-AUV comparison \\
        \texttt{C06\_representative} & 2023-12-22 & C06 & Multi-AUV comparison and IQR10 diagnosis \\
        \texttt{October\_control} & 2024-10-30 & C02 & Corrected October control case \\
        \hline
    \end{tabular}
\end{table}

The variables handled across the pipeline are TEMP or TEMPpred, STD, latitude, longitude, bathymetry and valid-mask information. The planner-ready NetCDF files embed the reward maps into the Lucrezia planner grid, while the ROI mask restricts valid cells to the operational region.

\section{Problem Formulation}

Let $\mathbf{s}$ denote a valid candidate sampling location in the ROI. The baseline sampling value is the normalised STD value $\sigma_{\mathrm{norm}}(\mathbf{s})$. The proposed enriched maps add a normalised descriptor $D_{\mathrm{norm}}(\mathbf{s})$ associated with the predicted regime class:

\begin{equation}
R^{\mathrm{base}}(\mathbf{s}) =
\sigma_{\mathrm{norm}}(\mathbf{s}).
\label{eq:ch4_baseline_reward}
\end{equation}

\begin{equation}
R^{\mathrm{enriched}}(\mathbf{s}) =
(1-\alpha)\,\sigma_{\mathrm{norm}}(\mathbf{s})
+
\alpha\,D_{\mathrm{norm}}(\mathbf{s}).
\label{eq:ch4_single_reward}
\end{equation}

For the multi-AUV vehicle-specific experiments, the two vehicles are assigned complementary region maps:

\begin{equation}
R^{(1)}(\mathbf{s}) =
w_{\sigma}\,\sigma_{\mathrm{norm}}(\mathbf{s})
+
w_A\,R_A(\mathbf{s}),
\label{eq:ch4_multi_reward_1}
\end{equation}

\begin{equation}
R^{(2)}(\mathbf{s}) =
w_{\sigma}\,\sigma_{\mathrm{norm}}(\mathbf{s})
+
w_B\,R_B(\mathbf{s}).
\label{eq:ch4_multi_reward_2}
\end{equation}

These equations define the final reward-map formulations used in this dissertation. The descriptor terms complement STD and are evaluated through their effect on planner outputs. Descriptor-only settings, such as $\alpha=1$ or $w_{\sigma}=0$, are treated as sensitivity tests rather than the main operational recommendation.

\section{Operational Constraints for Fair Comparison}

All comparisons are constrained so that changes in trajectory behaviour can be attributed to reward-map shaping rather than changes in the domain or planner configuration. The following constraints are used across the corrected experiments:

\begin{itemize}
    \item the same X490 ROI, grid and common valid mask are used for all reward maps;
    \item baseline and enriched maps are normalised over the valid ROI cells;
    \item the planner receives equivalent candidate and feasibility information for baseline and enriched cases;
    \item the mission-duration setting is fixed within each comparison, with the main corrected Step12A and Step12B analyses using 12 h missions;
    \item the Lucrezia planner itself is not modified; only the input reward maps are changed;
    \item planner reruns are analysed using their recorded manifests, solver diagnostics and generated trajectories.
\end{itemize}

\section{Evaluation Metrics}

The evaluation combines reward-based, geometric, diagnostic and computational metrics. The central uncertainty metric is collected STD, usually expressed together with STD retention relative to the STD-only baseline. Descriptor-related metrics include collected descriptor score, descriptor gain and distance-to-boundary quantities for boundary-distance descriptors. Trajectory geometry is assessed with route length, reward per distance, path difference from the baseline and qualitative route-map figures.

For multi-AUV experiments, fleet-level metrics include collected STD, region coverage, overlap and specialisation. Solver status and runtime are reported because operational feasibility is part of the planning problem. Finally, IQR10 is used as an external post-processing diagnostic of recent temporal variability. It is not part of the planner objective and it is not a direct data-assimilation metric.

\section{Chapter Summary}

This chapter defined the study case, the canonical data products and the controlled comparison problem. The planner remains uncertainty-driven through STD, while prototype-derived descriptors are introduced as complementary reward-map shaping terms. The following chapter describes how the offline regime library and the planning-day reward maps are constructed.


\chapter{Proposed Approach} \label{chap:proposed_approach}

This chapter presents the proposed methodology. The approach has two stages. The offline stage converts historical surface temperature fields into recurrent full-field regime classes, prototypes and prototype-derived descriptor maps. The planning-day stage assigns a daily TEMPpred map to one of those classes, retrieves the corresponding descriptors, combines them with STD and exports planner-ready reward maps.

\section{Methodological Overview}

The method is deliberately modular. The regime-inference stage is independent of the planner, and the planner is used without modifying its objective solver. The interface between both stages is the reward map. This design allows the thesis to evaluate whether descriptor-informed reward maps alter trajectory behaviour while keeping the operational uncertainty-driven baseline intact.

\begin{figure}[htbp]
    \centering
    \includegraphics[width=\linewidth]{figures/results/regime_inference_pipeline_diagram.pdf}
    \caption{Offline regime-inference pipeline. Historical surface temperature fields are cropped to the X490 ROI, transformed into compact representations, clustered into recurrent full-field regime classes and converted into prototype-derived descriptors for downstream reward-map construction.}
    \label{fig:regime_inference_pipeline}
\end{figure}

\begin{figure}[htbp]
    \centering
    \includegraphics[width=\linewidth]{figures/results/planner_inference_reward_generation_diagram.pdf}
    \caption{Planning-day inference and reward-map generation. TEMPpred selects the prototype class, prototype-derived descriptors are retrieved, STD remains the uncertainty baseline and the resulting reward maps are passed to the Lucrezia AUV planner. The output trajectories are then compared using STD, descriptor, geometry and diagnostic metrics.}
    \label{fig:planner_inference_reward_generation}
\end{figure}

\section{Offline Regime and Descriptor Library Construction}

The offline stage starts from the 370 historical surface temperature maps in the X490 ROI. After common-mask construction and normalisation, each daily field is represented using a Fossum-style compact model. The final canonical configuration uses patches of size $40\times24$, a dictionary size of 4, seed 11, StandardScaler enabled, Ward hierarchical clustering and an SD threshold of 0.25. This produces six recurrent temperature-regime classes.

In this dissertation, a regime class is not a cell-wise segmentation label. It is a recurrent full-field temperature pattern. Each class is represented by a prototype map obtained from its member days. The prototypes are then interpreted as homogeneous, single-gradient or multi-regime patterns. This interpretation determines which prototype-derived descriptors are meaningful for each class.

The descriptor family includes blended boundary scores, explicit boundary-distance maps in cells and kilometres, radius-based boundary-distance scores for radii 1, 2, 3, 5 and 8 cells, representative zones, cold and warm regions, and interest maps. For vehicle-specific multi-AUV experiments, complementary maps \texttt{region\_A} and \texttt{region\_B} are derived from the prototype regions.

\section{Planning-Day Class Assignment}

For a planning day $t$, the TEMPpred field is projected into the same regime-inference representation and assigned to one of the six previously discovered classes. If the assigned class is $C_k$, the planner retrieves the descriptor maps associated with the prototype of $C_k$. This step is important: descriptors used for planning are prototype-derived descriptors of the assigned class, not descriptors recalculated directly from STD.

The corrected Step11Y audit established the prototype-based lineage used in the final Step12 experiments. Earlier mixed-input logic was identified for some region-map fallback paths, especially for the October control case. The corrected workflow therefore uses the audited prototype-based arrays when generating planner inputs.

\section{Reward-Map Shaping}

Reward-map shaping combines the uncertainty baseline with one descriptor at a time in the single-AUV experiments. The baseline reward is

\begin{equation}
R^{\mathrm{base}}(\mathbf{s}) =
\sigma_{\mathrm{norm}}(\mathbf{s}).
\label{eq:ch5_baseline_reward}
\end{equation}

The single-AUV enriched reward is

\begin{equation}
R^{\mathrm{enriched}}(\mathbf{s}) =
(1-\alpha)\,\sigma_{\mathrm{norm}}(\mathbf{s})
+
\alpha\,D_{\mathrm{norm}}(\mathbf{s}),
\label{eq:ch5_single_reward}
\end{equation}

where $\alpha \in [0,1]$ controls the balance between STD and the selected descriptor. When $\alpha=0$, the route is the STD-only baseline. When $\alpha=1$, the route is descriptor-only and is interpreted as a sensitivity extreme.

For vehicle-specific multi-AUV experiments, the maps are

\begin{equation}
R^{(1)}(\mathbf{s}) =
w_{\sigma}\,\sigma_{\mathrm{norm}}(\mathbf{s})
+
w_A\,R_A(\mathbf{s}),
\label{eq:ch5_multi_reward_1}
\end{equation}

\begin{equation}
R^{(2)}(\mathbf{s}) =
w_{\sigma}\,\sigma_{\mathrm{norm}}(\mathbf{s})
+
w_B\,R_B(\mathbf{s}).
\label{eq:ch5_multi_reward_2}
\end{equation}

Here, $R_A$ and $R_B$ are complementary prototype-derived region maps. The weights are interpreted as reward-map shaping parameters, not as evidence that one vehicle directly reduces model uncertainty in the corresponding region.

\section{Planner Interface}

The reward maps are exported as planner-ready arrays and NetCDF inputs for the Lucrezia PCVRP/orienteering planner. The planner receives a spatial reward field, feasibility information and mission settings, then returns one or more trajectories. Since the planner is not modified, all methodological changes enter through the construction of the reward maps.

\section{Interpretation of Descriptor Weights}

Descriptor weights are interpreted conservatively. High descriptor weights can reveal whether a descriptor is capable of changing route geometry, but such runs may also sacrifice STD reward, increase runtime or produce operationally less robust trajectories. For this reason, the dissertation gives preference to configurations that retain a substantial fraction of STD while adding interpretable descriptor coverage. Descriptor-only experiments are useful sensitivity tests, but they are not presented as the primary operational solution.

\section{Chapter Summary}

The proposed approach converts historical temperature structure into a prototype descriptor library and uses planning-day TEMPpred maps only to select the relevant class. STD remains the uncertainty baseline. The next chapter documents how this methodology was implemented and audited.


\chapter{Implementation} \label{chap:implementation}

This chapter documents the executed implementation pipeline and the audit artefacts used to support reproducibility. The emphasis is on the completed workflow rather than on intended future processing.

\section{Executed Pipeline}

The implementation is organised into numbered processing steps. Each step writes explicit outputs, manifests or reports. Table~\ref{tab:implementation_steps} summarises the artefacts used in this dissertation.

\begin{table}[htbp]
    \centering
    \small
    \caption{Executed implementation steps and main artefacts used in the thesis.}
    \label{tab:implementation_steps}
    \begin{tabular}{p{0.16\linewidth}p{0.32\linewidth}p{0.38\linewidth}}
        \hline
        Step & Purpose & Main artefacts \\
        \hline
        Step00 & Canonical ROI and dataset preparation & Temperature stack, mask, coordinates, bathymetry, dates and dataset metadata \\
        Step05 & Canonical compact regime model & Sparse codes, dictionary learning outputs, dendrogram, class assignments and prototypes \\
        Step08 & Final descriptor generation & Descriptor maps, descriptor definitions, quality flags and prototype summaries \\
        Step09B & TEMPpred class assignment & Planning-day class assignment for representative and control dates \\
        Step11Y/Z & Planner-input audit and correction & Prototype-based reward arrays, lineage checks and corrected planner input maps \\
        Step12A & Single-AUV sensitivity & Planner inputs, runs, manifests, metrics, solver diagnostics and figures \\
        Step12B & Multi-AUV sensitivity & Fleet and vehicle metrics, manifests, solver diagnostics and figures \\
        Diagnostics & Interpretation support & Descriptor-weight statistics, IQR10 maps, IQR10 route metrics and trade-off figures \\
        \hline
    \end{tabular}
\end{table}

\section{Step00: Canonical ROI and Dataset Preparation}

Step00 prepares the canonical X490 surface temperature dataset. The audited folder contains \texttt{X\_surface\_370\_roi\_x490.npy}, normalised temperature arrays, the common valid mask, bathymetry, latitude, longitude, projected coordinates and the dates CSV. The metadata verifies the date range from 2023-10-28 to 2024-10-31, the $370 \times 72 \times 117$ shape and the 8004-cell common mask.

The same Step00 outputs are used by the compact-model steps, descriptor generation and diagnostic post-processing. This avoids changes in ROI definition between the regime model and the planner evaluation.

\section{Step01--Step05: Regime Discovery and Canonical Configuration}

The canonical compact representation uses the audited Step05 folder. The configuration is patch size $40\times24$, dictionary size 4, seed 11, StandardScaler enabled, Ward hierarchical clustering and SD threshold 0.25. The final feature matrix has 370 rows, one per day, and the clustering produces six classes with sizes 41, 70, 50, 107, 30 and 72 days.

The Step05 artefacts include class prototypes, dendrograms, class-size plots, sparse codes and canonical reports. These artefacts are used both as evidence of reproducibility and as inputs to descriptor generation.

\section{Step08: Prototype-Derived Descriptor Generation}

Step08 converts the six class prototypes into descriptor maps. The final descriptor set includes blended boundary scores, explicit boundary-distance maps, radius-based boundary-distance scores, representative zones, cold and warm regions and interest maps. All planner maps are normalised over the Step00 valid mask. For homogeneous classes, boundary-oriented maps can be zero or weak because no internal boundary is supported by the prototype interpretation.

\begin{figure}[htbp]
    \centering
    \includegraphics[width=\linewidth]{figures/results/step08_descriptor_maps_by_class_panel_shared_prototype_scale.png}
    \caption{Prototype-derived descriptors generated in Step08. The descriptor maps are associated with full-field prototype classes and are later retrieved by planning-day class assignment.}
    \label{fig:step08_descriptor_maps}
\end{figure}

\section{Step09B and Step11Y/Z: Class Assignment and Planner-Input Audit}

The planning-day TEMPpred maps are assigned to the canonical classes before descriptor retrieval. The final audited cases assign C01\_representative to C01, C06\_representative to C06 and October\_control to C02. The Step11Y audit identified mixed-input behaviour in older region-map fallback logic and produced corrected \texttt{prototype\_based\_*} arrays, including the boundary-distance score variants used in Step12A.

\begin{figure}[htbp]
    \centering
    \includegraphics[width=0.95\linewidth]{figures/results/step11y_planner_input_logic_diagram.png}
    \caption{Planner-input logic audited in Step11Y. The corrected workflow uses planning-day TEMPpred only for class assignment and uses prototype-based descriptor maps for reward-map construction.}
    \label{fig:step11y_planner_input_logic}
\end{figure}

\section{Step12A: Single-AUV Experiments}

Step12A evaluates one AUV over the C01 representative case with a 12 h mission duration. The tested descriptors are the blended boundary score, boundary-distance scores at r1, r3 and r5, and the interest map. The alpha sweep uses five values, including the STD-only baseline $\alpha=0$ and the descriptor-only sensitivity extreme $\alpha=1$.

The output folder contains planner inputs, planner run directories, the run manifest, the main metrics CSV, solver diagnostics, runtime summaries, IQR10 post-processing and figures. The final Step12A folder used here is the corrected 20260608 run, which uses prototype-based maps only and does not use TEMPpred as the objective.

\section{Step12B: Multi-AUV Vehicle-Specific Experiments}

Step12B evaluates two-AUV configurations for C01\_representative, C06\_representative and October\_control. The baseline uses a shared STD map. Vehicle-specific configurations combine STD with complementary prototype region maps so that AUV1 and AUV2 can be encouraged towards different parts of the assigned prototype. The implementation is a wrapper/proxy experiment: individual vehicle reward maps are prepared and solved, then combined into fleet-level metrics. The Lucrezia planner objective is not modified.

\section{IQR10 and Descriptor Diagnostics}

The IQR10 diagnostic is computed after planner execution. For each planning date, the interquartile range of temperature over the previous 10 available days is mapped over the ROI:

\begin{equation}
\mathrm{IQR10}(\mathbf{s},t)
=
P_{75}\{T(\mathbf{s},\tau):\tau<t,\tau\in W_{10}(t)\}
-
P_{25}\{T(\mathbf{s},\tau):\tau<t,\tau\in W_{10}(t)\}.
\label{eq:iqr10_definition}
\end{equation}

The IQR10 maps are then sampled along the existing trajectories to compute collected IQR10, IQR10 per kilometre and gain relative to the baseline. This is post-processing only and does not rerun the planner.

\section{Reproducibility and Lineage}

The implementation stores reports, CSV summaries, JSON checks, NetCDF planner inputs, NPZ arrays and figures for each major step. When multiple output folders existed, the final writing prioritised the newest corrected or audited artefacts: Step11Y 20260606 for prototype-based input correction, Step12A 20260608 for single-AUV results and Step12B 20260608 for multi-AUV results. Older folders are treated as historical lineage or debugging context.

\section{Chapter Summary}

This chapter documented the executed pipeline and the artefacts used for reproducibility. The next chapter defines the verification, validation and experimental design used to interpret the outputs.


\chapter{Verification, Validation and Experimental Design} \label{chap:validation}

This chapter describes how the pipeline and planner experiments were checked and how the results are interpreted. The goal is not to validate an assimilation system, but to evaluate whether descriptor-enriched reward maps change planner behaviour in a controlled and interpretable way.

\section{Data and Grid Validation}

The first validation layer checks that temperature, mask, coordinates, bathymetry and dates are consistent. The Step00 metadata verifies the number of maps, date range, ROI shape, common-mask size, coordinate system and generated PNG count. The same mask is used to normalise descriptor maps and to define valid planner cells.

\section{Prototype and Descriptor Validation}

The compact-model validation checks that the final canonical configuration produces six recurrent full-field temperature classes and that the class sizes are plausible. Prototype interpretation then separates homogeneous, single-gradient and multi-regime structures. Descriptor validation checks whether descriptors are finite on the valid mask, whether boundary-oriented descriptors are meaningful only where boundaries exist and whether descriptor maps can be propagated into planner-ready arrays.

\begin{figure}[htbp]
    \centering
    \includegraphics[width=0.95\linewidth]{figures/results/canonical_class_prototypes_panel.png}
    \caption{Canonical six-class prototype panel from Step05. The classes represent recurrent full-field temperature patterns rather than cell-wise labels.}
    \label{fig:canonical_prototypes}
\end{figure}

\section{TEMPpred Class Assignment Validation}

For each planning day, TEMPpred is assigned to a previously discovered class. The assigned class determines which prototype descriptors are used. The Step11Y audit then checks whether planner maps are actually prototype-based and whether older fallback logic introduced mixed inputs. This validation is essential because a visually plausible route comparison can still be methodologically inconsistent if its descriptor maps come from the wrong source.

\section{Single-AUV Experimental Design}

The single-AUV experiment uses the C01 representative case, a 12 h mission duration and a descriptor-alpha sweep. The descriptors are the blended boundary score, boundary-distance scores at r1, r3 and r5, and the interest map. For each descriptor, $\alpha=0$ provides the STD-only baseline and larger $\alpha$ values increase the descriptor contribution.

The main metrics are collected STD, STD retention, collected descriptor score, route length, reward efficiency, path difference from baseline, runtime and solver status. Boundary-distance variants are interpreted as a sensitivity test of descriptor spatial scale: r1 is narrow, r3 is intermediate and r5 is broader.

\section{Multi-AUV Experimental Design}

The multi-AUV experiments compare a shared-STD baseline with vehicle-specific region-map reward shaping. The corrected Step12B runs use C01\_representative, C06\_representative and October\_control, all with 12 h missions. Fleet-level metrics include collected STD, STD retention, region coverage, route overlap, specialisation and runtime. These metrics assess whether complementary region maps can reduce redundant behaviour or encourage vehicles to sample different regime-related areas.

\section{IQR10 as an external diagnostic}

IQR10 measures recent temporal variability using the previous 10 available days of surface temperature. It is not part of the reward function, is not passed to the planner and is not a direct assimilation metric. It also does not prove forecast improvement or uncertainty reduction. Its purpose is diagnostic: after trajectories have been generated, IQR10 indicates whether a route sampled regions that had recently variable temperature values and may therefore be relevant for future assimilation-aware studies.

\section{Solver and Runtime Validation}

Every planner run is interpreted together with its solver status and runtime. Successful runs are used for quantitative recommendations. Timeouts or failed runs are not silently discarded; they are used to identify operationally fragile configurations. This is particularly important for descriptor-only or very narrow descriptor settings, where a map can become too sparse or computationally difficult.

\section{Interpretation Rules}

The interpretation follows four rules. First, STD-only remains the operational baseline. Second, an enriched route is useful only if it preserves an acceptable amount of STD while adding interpretable descriptor sampling or reducing redundancy. Third, descriptor-only results are sensitivity tests and are not automatically operational recommendations. Fourth, IQR10 is used only as supporting evidence about recent variability, not as proof of assimilation benefit.

\section{Chapter Summary}

The verification and validation design links each result to its data lineage, descriptor source, planner status and diagnostic interpretation. The next chapter reports the results using this framework.


\chapter{Results and Discussion} \label{chap:results}

This chapter presents the corrected results and discusses what they indicate about regime-aware reward-map shaping. The emphasis is on controlled planner behaviour, not on direct data-assimilation improvement.

\section{Canonical Dataset Validation}

The Step00 audit confirms that the canonical X490 dataset is internally consistent. The 370 surface temperature maps span 2023-10-28 to 2024-10-31, the cropped stack has shape $370 \times 72 \times 117$ and the common valid mask contains 8004 cells. The generated daily PNG set is complete, and the metadata confirms that the corrected X490 ROI was applied.

This matters because every later step depends on the same spatial support. The regime model, descriptor maps, planning-day maps and IQR10 diagnostics are all interpreted over the same ROI and mask.

\section{Regime Discovery Results}

The canonical compact-model configuration produced six recurrent full-field temperature classes. Table~\ref{tab:class_sizes} reports the class sizes from Step05.

\begin{table}[htbp]
    \centering
    \small
    \caption{Class sizes obtained with the canonical compact-regime configuration.}
    \label{tab:class_sizes}
    \begin{tabular}{p{0.18\linewidth}p{0.22\linewidth}p{0.28\linewidth}}
        \hline
        Class & Number of days & Fraction of 370 days \\
        \hline
        C01 & 41 & 11.08\% \\
        C02 & 70 & 18.92\% \\
        C03 & 50 & 13.51\% \\
        C04 & 107 & 28.92\% \\
        C05 & 30 & 8.11\% \\
        C06 & 72 & 19.46\% \\
        \hline
    \end{tabular}
\end{table}

The prototypes are interpretable as spatial temperature patterns. C01 and C06 are the main multi-regime or boundary-rich classes in the final interpretation, while C03 and C04 are more homogeneous. This distinction is important because boundary-oriented descriptors are meaningful only when the prototype supports a boundary or strong transition.

\section{Descriptor Library Results}

Step08 generated the prototype-derived descriptor library and exported normalised planner maps. The descriptor ranking reported by Step08 indicates that C01, C05 and C06 are prominent for boundary and gradient descriptors, while interest-map ranking emphasises C01, C06 and C05. The descriptor-weight diagnostic further supports a cautious interpretation: C01 has very high prototype heterogeneity, C06 has medium heterogeneity and the homogeneous classes have low heterogeneity.

The boundary-distance descriptors separate pure proximity to a boundary from the older blended \texttt{boundary\_score}. This separation is useful because it allows spatial scale to be tested explicitly. The r1 descriptor is narrow and emphasises cells immediately adjacent to the boundary; r3 provides an intermediate band; r5 is broader and can spread reward over a larger surrounding region.

\section{Prototype-Based Planner Input Audit}

The corrected Step11Y audit confirmed that the Step10F boundary inputs matched Step08 prototype descriptors for the three evaluated cases. It also identified mixed-input evidence in older region-map fallback paths for October\_control. The final Step12 experiments therefore use the corrected \texttt{prototype\_based\_*} arrays. This lineage correction is a central methodological result: the final planner comparison is based on the assigned prototype descriptors rather than on direct TEMPpred-derived region fallbacks.

\section{Single-AUV Results}

The corrected Step12A run evaluates the C01 representative case for a 12 h mission. It contains 21 physical planner runs and 25 logical sensitivity rows. The report confirms that prototype-based maps were used and that TEMPpred was not used as the objective.

\begin{figure}[htbp]
    \centering
    \includegraphics[width=\linewidth]{figures/results/step12a_C01_representative_12h_boundary_distance_descriptor_comparison_alpha050.png}
    \caption{Single-AUV route comparison for the C01 representative case using boundary-distance descriptor variants. The figure illustrates how descriptor scale changes the selected trajectory while the same planning domain and baseline STD context are retained.}
    \label{fig:step12a_boundary_distance_comparison}
\end{figure}

The STD-only baseline collected 90.650 STD units with a route length of 38.053 km in the reported metrics. Table~\ref{tab:step12a_selected_results} summarises selected successful runs and the recommended configurations reported by Step12A.

\begin{table}[htbp]
    \centering
    \small
    \caption{Selected Step12A single-AUV results for the C01 representative 12 h case.}
    \label{tab:step12a_selected_results}
    \begin{tabular}{p{0.32\linewidth}p{0.12\linewidth}p{0.16\linewidth}p{0.16\linewidth}p{0.13\linewidth}}
        \hline
        Descriptor & $\alpha$ & Collected STD & STD retention & Runtime (s) \\
        \hline
        STD-only baseline & 0.00 & 90.650 & 1.000 & 1524.010 \\
        Boundary score & 0.50 & 89.061 & 0.982 & 924.108 \\
        Boundary-distance r1 & 0.75 & 84.847 & 0.936 & 680.524 \\
        Boundary-distance r3 & 1.00 & 94.488 & 1.042 & 214.716 \\
        Boundary-distance r5 & 1.00 & 93.691 & 1.034 & 287.738 \\
        Interest map & 0.50 & 94.736 & 1.045 & 1344.849 \\
        \hline
    \end{tabular}
\end{table}

The results indicate that descriptor-enriched routes can substantially change geometry while preserving or, in some cases, slightly exceeding the baseline collected STD. However, this does not mean that descriptor-only planning is generally preferable. The r3 and r5 descriptor-only cases scored well for this specific C01 setup, but descriptor-only settings remain sensitivity extremes. The \texttt{boundary\_score} at $\alpha=0.50$ and \texttt{interest\_map} at $\alpha=0.50$ are more conservative examples because they preserve STD while keeping a direct uncertainty contribution in the reward map.

\begin{figure}[htbp]
    \centering
    \includegraphics[width=0.85\linewidth]{figures/results/step12a_STD_vs_descriptor_tradeoff.png}
    \caption{Step12A trade-off between collected STD and descriptor-related behaviour. The figure supports the interpretation that useful enriched routes should be judged by STD retention, descriptor response and operational feasibility together.}
    \label{fig:step12a_tradeoff}
\end{figure}

The r1 descriptor behaved as a narrow boundary-distance map. At intermediate weights it produced two timeouts, and at $\alpha=1.00$ it collected substantially less STD than the baseline. This supports the interpretation that r1 can be too spatially restrictive. The r5 descriptor is broader and less fragile, but it may dilute the precise boundary focus. The r3 descriptor provides a useful compromise in this case because it preserves strong STD performance while still expressing boundary proximity.

\section{Multi-AUV Results}

The corrected Step12B run evaluates three cases with 12 h missions. It contains 27 fleet rows and 54 vehicle rows. The report confirms that prototype-based maps were used, TEMPpred was not used as the objective and no overlap/proximity penalty was enabled in the final comparison.

\begin{figure}[htbp]
    \centering
    \includegraphics[width=0.95\linewidth]{figures/results/step12b_C01_representative_12h_best_configuration_panel.png}
    \caption{Best C01 representative multi-AUV configuration from Step12B. Vehicle-specific prototype-region maps can encourage route specialisation while retaining a shared uncertainty-driven component.}
    \label{fig:step12b_c01_best}
\end{figure}

Table~\ref{tab:step12b_recommendations} reports the best recommendation for each case in the Step12B summary.

\begin{table}[htbp]
    \centering
    \small
    \caption{Best Step12B vehicle-specific recommendations by case.}
    \label{tab:step12b_recommendations}
    \begin{tabular}{p{0.25\linewidth}p{0.20\linewidth}p{0.14\linewidth}p{0.14\linewidth}p{0.14\linewidth}}
        \hline
        Case & Strategy & STD retention & Specialisation & Overlap \\
        \hline
        C01\_representative & 70/30 STD-region & 0.891 & 0.452 & 0.003 \\
        C06\_representative & 60/40 STD-region & 0.856 & 0.551 & 0.003 \\
        October\_control & 25/75 role swap & 0.974 & 0.000 & 1.000 \\
        \hline
    \end{tabular}
\end{table}

The C01 and C06 results suggest that vehicle-specific maps can encourage specialisation with very low overlap. The October\_control result is less favourable: the recommended role-swap strategy retained STD but had complete overlap and no specialisation in the reported metric. This supports the broader conclusion that multi-AUV benefits are case-dependent and require stronger coordination mechanisms if route redundancy is to be controlled reliably.

The all-descriptor vehicle-specific 0/100 setting is not robust. In Step12B it had a lower mean STD retention and a success rate of 0.667 in the duration summary. The solver diagnostics recorded two failed runs for October\_control under \texttt{vehicle\_specific\_00100}, where the reward input had zero maximum value. This reinforces the preference for mixed STD-descriptor weights.

\section{IQR10 Results}

IQR10 was computed as a post-processing diagnostic for the existing trajectories. The windows used 10 previous days for all evaluated cases. The validation values are shown in Table~\ref{tab:iqr10_windows}.

\begin{table}[htbp]
    \centering
    \small
    \caption{IQR10 window validation for the corrected Step12 diagnostics.}
    \label{tab:iqr10_windows}
    \begin{tabular}{p{0.24\linewidth}p{0.16\linewidth}p{0.22\linewidth}p{0.12\linewidth}p{0.12\linewidth}}
        \hline
        Case & Date & Previous-day window & Mean IQR10 & Finite cells \\
        \hline
        C01\_representative & 2024-08-24 & 2024-08-14 to 2024-08-23 & 0.7699 & 8004 \\
        C06\_representative & 2023-12-22 & 2023-12-12 to 2023-12-21 & 0.2173 & 8004 \\
        October\_control & 2024-10-30 & 2024-10-20 to 2024-10-29 & 1.0293 & 8004 \\
        \hline
    \end{tabular}
\end{table}

For Step12A, the top routes by collected IQR10 were the STD-only baseline rows, with collected IQR10 of 95.298 and 32.03\% of the path in the top 10\% IQR10 cells. Thus, in the C01 single-AUV case, enriched routes did not improve the IQR10 diagnostic relative to the baseline.

For Step12B, the IQR10 results were mixed across days. The C06 representative case showed positive IQR10 gains for the 90/10 and 80/20 vehicle-specific strategies. The best reported gain was 10.066 IQR10 units, or 6.83\%, for \texttt{vehicle\_specific\_9010}, with STD retention 1.076. In contrast, October\_control showed negative IQR10 gain for the reported enriched strategies relative to the shared-STD baseline. These results support the interpretation that IQR10 is scenario-dependent and should remain a diagnostic rather than a design claim.

\begin{figure}[htbp]
    \centering
    \includegraphics[width=0.95\linewidth]{figures/results/step12b_iqr10_maps_three_days_normalized.png}
    \caption{Normalised IQR10 maps for the three evaluated planning dates. IQR10 indicates recent temporal variability over the previous 10 available days and is used only as a post-processing diagnostic.}
    \label{fig:step12b_iqr10_three_days}
\end{figure}

\begin{figure}[htbp]
    \centering
    \includegraphics[width=0.85\linewidth]{figures/results/step12a_iqr10_map_and_selected_routes.png}
    \caption{Step12A IQR10 map and selected routes for the C01 representative case. The diagnostic overlays existing trajectories on recent temporal variability but does not modify the planner reward.}
    \label{fig:step12a_iqr10_routes}
\end{figure}

\section{Operational Feasibility}

Runtime and solver status affect whether a reward-map strategy is operationally plausible. In Step12A, 19 runs succeeded and 2 runs timed out. Both timeouts occurred for the r1 boundary-distance descriptor at $\alpha=0.25$ and $\alpha=0.50$. In Step12B, 29 runs succeeded, 20 were reused from existing runs and 2 failed. The failed cases occurred for October\_control under descriptor-only vehicle-specific maps with zero-valued reward inputs.

These diagnostics show that map validity alone is not enough. A descriptor can be mathematically well-defined but still produce a sparse, weak or operationally fragile reward map. This is another reason to favour conservative mixed weights.

\section{Discussion}

The results support four main conclusions. First, the corrected pipeline can translate historical temperature-regime structure into planner-compatible descriptors. Second, descriptor-enriched rewards can change route geometry and, in selected cases, preserve STD reward. Third, the effect is case-dependent: C01, C06 and October\_control do not respond in the same way. Fourth, conservative descriptor weights are easier to justify than descriptor-only formulations because they preserve the uncertainty-driven role of STD while allowing descriptors to add spatial structure.

The comparison between boundary descriptors is particularly informative. The old boundary score remains useful as a blended front descriptor, but the pure boundary-distance variants expose the effect of spatial scale more clearly. The r1 map is narrow and can be operationally fragile; the r5 map is broader and less restrictive; r3 offers an intermediate compromise for C01. The interest map performs well in the C01 single-AUV summary at $\alpha=0.50$, but it should still be treated as a complementary descriptor rather than as a replacement for STD.

The multi-AUV results indicate that vehicle-specific maps can encourage specialisation in C01 and C06, but overlap and role assignment remain limitations. In particular, October\_control shows that assigning complementary maps does not guarantee complementary trajectories. Future multi-AUV work should therefore include explicit overlap penalties, joint optimisation or communication-aware coordination.

\section{Chapter Summary}

The corrected results show that regime-aware descriptors can complement an uncertainty-driven planner, but they do not universally dominate the STD-only baseline. The most defensible interpretation is that descriptors provide additional spatial context for planning, especially in boundary-rich regimes, while STD remains a strong operational baseline. No direct data-assimilation improvement or forecast-skill improvement is claimed.


\chapter{Conclusions} \label{chap:conclusions}

This chapter summarises the thesis findings, contributions, limitations and future research directions.

\section{Main Findings}

The dissertation developed and evaluated a reproducible pipeline for converting temperature-regime information into planner-compatible descriptors. Historical surface temperature maps were used to identify recurrent full-field regimes, generate class prototypes and derive descriptors such as boundary scores, boundary-distance maps, representative zones, cold and warm regions and interest maps.

The planner remains uncertainty-driven. STD is the baseline reward, and descriptor maps complement STD through controlled reward-map shaping. The single-AUV results show that descriptors can change trajectory geometry and sometimes preserve high STD reward, but the effect depends on descriptor type and weight. The boundary-distance analysis suggests that r1 can be too narrow, r5 can be broad and r3 can provide a useful intermediate scale in the C01 representative case.

The multi-AUV results show that vehicle-specific descriptor maps can encourage specialisation in some cases, especially C01 and C06, but redundancy remains a limitation. October\_control demonstrated that complementary input maps do not automatically produce complementary trajectories.

IQR10 provided an external diagnostic of recent temporal variability. It helped interpret whether existing trajectories sampled recently variable regions, but it was not part of the reward function and does not measure assimilation performance.

\section{Contributions}

The main contributions are:

\begin{itemize}
    \item a reproducible X490 ROI processing chain for 370 historical surface temperature maps;
    \item a canonical compact-model regime library with six recurrent full-field temperature classes;
    \item a prototype-derived descriptor library suitable for reward-map shaping;
    \item a corrected planning-day workflow in which TEMPpred selects the class and STD remains the baseline reward;
    \item single-AUV and multi-AUV planner evaluations using prototype-based descriptors;
    \item post-processing diagnostics for STD retention, descriptor response, multi-AUV overlap and IQR10 temporal variability.
\end{itemize}

\section{Limitations}

The results are case-dependent and are based on a limited set of evaluated planning dates. The experiments evaluate planner behaviour rather than direct forecast improvement. Data assimilation is not implemented in this thesis, so the results should not be interpreted as evidence of uncertainty reduction in a model state.

The multi-AUV experiments use a wrapper/proxy strategy for vehicle-specific maps and do not implement a fully joint anti-overlap objective. Some descriptor-only configurations were operationally fragile, including timeout and failed-run cases. Finally, the descriptor interpretation depends on the quality of the prototype classes and on the representativeness of the historical temperature archive.

\section{Future Work}

Future work should connect the planner outputs to direct data-assimilation experiments so that observation impact can be measured explicitly. Online replanning could update class assignment and reward maps as new observations arrive. Multi-AUV coordination should be extended with explicit overlap penalties, joint optimisation and communication-aware constraints. The evaluation should also be expanded across more dates, seasons and operational conditions, followed by validation with real AUV deployments.

\section{Final Remarks}

The thesis shows that temperature-derived regime descriptors can provide useful complementary spatial information for an uncertainty-driven AUV planner. The strongest conclusion is not that descriptors replace STD, but that they can support more interpretable reward-map sensitivity analyses and reveal how trajectories respond to prototype-based regime information. This is a step towards assimilation-aware mission design, while leaving direct assimilation impact as a necessary future validation stage.
